\documentclass[10pt,a4paper]{amsart}
\usepackage[margin=19mm]{geometry}
\usepackage{amsmath,amssymb,mathtools}
\usepackage[scr=rsfs,cal=euler]{mathalfa}
\usepackage{xfrac}
\usepackage[unicode,hidelinks,bookmarks=false]{hyperref}
\newcommand{\ie}{i.e.\ }
\newcommand{\cf}{cf.\ }
\newcommand{\resp}{respectively\ }
\newcommand{\Subsection}{\subsection}
\providecommand{\nobreakdash}{\nobreak\hbox{-}}
\setlength{\emergencystretch}{2em}
\multlinegap0pt
\usepackage{bm}
\usepackage{bbm}
\usepackage{dsfont}
\usepackage{xspace}
\usepackage{cancel}
\usepackage{mathtools}
\usepackage{amsthm}
\makeatletter
\@ifundefined{theorem}{\newtheorem{theorem}{Theorem}}{}
\makeatother
\makeatletter
\def \Id{\mathbf{I}}
\def\bx{{\bf x}}
\def\by{{\bf y}}
\expandafter\ifx\csname equations\endcsname\relax
\newenvironment{equations}{\equation\aligned}{\endaligned\endequation}
\fi
\def \erre{\mathbb{R}}
\renewcommand{\le}{\leqslant}
\makeatother
\title[From Kinetic Theory to AI: a Rediscovery of High-Dimensional]{From Kinetic Theory to AI: a Rediscovery of High-Dimensional Divergences and Their Properties}
\author{Gennaro Auricchio, Giovanni Brigati, Paolo Giudici, Giuseppe Toscani}
\date{Source: arXiv:2507.11387v1 (CC BY 4.0)}
\begin{document}
\maketitle

\noindent\emph{Source and licence.} From Kinetic Theory to AI: a Rediscovery of High-Dimensional Divergences and Their Properties. Version arXiv:2507.11387v1 (CC BY 4.0), distributed under the Creative Commons Attribution 4.0 International licence, as recorded in the arXiv licence metadata for this exact version and shown on the abstract page https://arxiv.org/abs/2507.11387v1. Full rights notice: licenses/arxiv-2507.11387-CC-BY-4.0.txt.\par\smallskip

\noindent\textit{Editorial scope.} Counted unit: the Theorem whose source label is \texttt{None} in the article (source0.tex lines 2548--2557, source label \texttt{None}), delivered together with the article's own complete original proof of it (source0.tex lines 2559--2604, one \texttt{proof} environment of 46 lines). No mathematics is written, repaired or re-derived: statement and proof are the article's own text line for line. Required context retained uncounted: none. Every definition, display and label of those lines is retained unchanged. Omitted from this package and not redistributed: 1--2547, 2558--2558, 2605--4154. Nothing inside a retained excerpt is omitted or altered: every retained body is delivered exactly as the article's lines are written, so no omission receipt of any kind accompanies this package. Citations: the retained excerpts carry no citation command at all, so no bibliography entry of the article is transcribed and no reference is redistributed. Reference adaptation: none was needed and none was made; every reference inside the delivered unit resolves inside it, so no unresolved reference or citation is produced. Printed slips kept literal: no printed word, symbol, hypothesis or number of the counted statement or of its proof was altered. Layout and class: the article's own class is \texttt{article} with packages including graphicx, amsmath, amssymb, geometry, adjustbox, babel, fontenc, url, xcolor, tikz, hyperref, booktabs, enumitem, verbatim; the wrapper is the lane's pinned \texttt{amsart} editorial wrapper at 10pt on A4, which restates the theorem-like environments the retained text needs and adds the article's own macro definitions that the retained text needs (\texttt{\textbackslash Id}, \texttt{\textbackslash bx}, \texttt{\textbackslash by}, \texttt{\textbackslash erre}, \texttt{\textbackslash le}), with extra wrapper packages (bm, bbm, dsfont, xspace, cancel, mathtools). The delivered numbering is the unit's own. Assets: no retained span contains a figure environment, a graphics-inclusion command, a PDF-page inclusion, a table or a TikZ picture; no image file is copied into this package, the package declares no asset dependency and no image file is redistributed with it. Licence: version arXiv:2507.11387v1 of this article is distributed under the Creative Commons Attribution 4.0 International licence, as recorded in the arXiv licence metadata for this exact version and shown on the abstract page https://arxiv.org/abs/2507.11387v1. A full rights notice is delivered at licenses/arxiv-2507.11387-CC-BY-4.0.txt. Characters: every retained excerpt is pure ASCII, so the delivered engine, XeLaTeX with its default Latin Modern text font, reports no missing character.
\begin{theorem}
    Given two absolutely continuous probability distributions, namely $\mu$ and $\nu$, we have that
    \begin{equation}
        \frac{1}{2}\mathcal{E}^2_1(\mu,\nu)\le W_1(\mu,\nu)
    \end{equation}
    and
    \begin{equation}
        \frac{\Gamma\left(\frac{n+1}2\right)}{2\pi^{(n+1)/2}}\mathcal{F}^2_{n+1}(\mu,\nu)\le W_1(\mu,\nu).
    \end{equation}
\end{theorem}

\begin{proof}
    Since $\mu$ and $\nu$ are absolutely continuous distributions, there exists two functions, namely $T$ and $S$ such that the following conditions hold
    \begin{itemize}
        \item $T_\#\mu=\nu$ and $S_\#\nu=\mu$;
        \item $T\circ S= S \circ T= \Id$, where $\Id$ is the identity function on $\erre^n$; and
        \item $W_1(\mu,\nu)=\int_{\erre^{n}}|\bx-T(\bx)|d\mu=\int_{\erre^{n}}|\by-S(\by)|d\nu$.
    \end{itemize}
    In particular, $T$ is the optimal transportation map from $\mu$ to $\nu$ and $S$ is the optimal transportation map from $\nu$ to $\mu$.
    % 

    % 
    We then have the following 
    \begin{align}
    \label{eq:ineqT}
        \nonumber\int_{\erre^n\times\erre^n}&|\bx-\by|\mu(d\bx)\nu(d\by)=\int_{\erre^n\times\erre^n}|\bx-T(\bx)+T(\bx)-\by|\mu(d\bx)\nu(d\by)\\
        \nonumber&\le \int_{\erre^n\times\erre^n}\big(|\bx-T(\bx)|+|T(\bx)-\by|\big)\mu(d\bx)\nu(d\by)\\
        \nonumber&=\int_{\erre^n}\int_{\erre^n}|\bx-T(\bx)|\mu(d\bx)\nu(d\by)+\int_{\erre^n}\int_{\erre^n}|T(\bx)-\by|\mu(d\bx)\nu(d\by)\\
        &=\int_{\erre^n}|\bx-T(\bx)|\mu(d\bx)+\int_{\erre^n}\int_{\erre^n}|\by'-\by|\nu(d\by')\nu(d\by)
    \end{align}
    since, owing to the fact that $T_\#\mu=\nu$, we have that 
    \begin{equation}
        % \label{eq:ineqT}
        \int_{\erre^n}\int_{\erre^n}|T(\bx)-\by|\mu(d\bx)\nu(d\by)=\int_{\erre^n}\int_{\erre^n}|\by'-\by|\nu(d\by')\nu(d\by).
    \end{equation}
    Similarly, by using $S$ instead of $T$, we have that
    \begin{equation}
        \label{eq:ineqS}
        \int_{\erre^n\times\erre^n}|\bx-\by|\mu(d\bx)\nu(d\by)\le \int_{\erre^n}|\by-S(\by)|\nu(d\by)+\int_{\erre^n}\int_{\erre^n}|\bx'-\bx|\mu(d\bx')\mu(d\bx),    
    \end{equation}
    % 
    Since $W_1(\mu,\nu)=\int_{\erre^n}|\by-S(\by)|\nu(d\by)$, summing up the inequalities \eqref{eq:ineqT} and \eqref{eq:ineqS}, we have
    \begin{align}
        \label{eq:ineq_tot}
        \nonumber2\int_{\erre^n\times\erre^n}|\bx-\by|\mu(d\bx)\nu(d\by)&\le 2W_1(\mu,\nu)+\int_{\erre^n}\int_{\erre^n}|\bx'-\bx|\mu(d\bx')\mu(d\bx)\\
        &\quad+ \int_{\erre^n}\int_{\erre^n}|\by'-\by|\nu(d\by')\nu(d\by),
    \end{align}
    which, by rearranging the terms and recalling the definition of the Energy Distance, we get
    \begin{align}
        \mathcal{E}_1^2(\mu,\nu)\le 2W_1(\mu,\nu),
    \end{align}
    hence
    \begin{equation}
        \frac{\Gamma\left(\frac{n+1}2\right)}{2\pi^{(n+1)/2}}\mathcal{F}^2_{n+1}(\mu,\nu)\le W_1(\mu,\nu)
    \end{equation}
    since $\Gamma(\frac{1}{2})=\sqrt{\pi}$.
\end{proof}

\end{document}
